\documentclass[11pt,a4paper]{article}

% --------------------
% Packages
% --------------------
\usepackage[margin=0.6in]{geometry}
\usepackage{amsmath, amssymb}
\usepackage{amsthm}
\usepackage{graphicx}
\usepackage{booktabs}
\usepackage{hyperref}
\usepackage{setspace}

\newtheorem{theorem}{Theorem}

% --------------------
% Document
% --------------------
\begin{document}


% --------------------
% Theorem
% --------------------
\begin{theorem}
Let $\mathbf{x} := (\mathbf{y}, \mathbf{z})$ and define
\[
\mathbf{F}(\mathbf{x}) :=
\begin{pmatrix}
\mathbf{G}(\mathbf{y}, \mathbf{z}) \\
\mathbf{H}(\mathbf{y}, \mathbf{z})
\end{pmatrix},
\qquad
\mathbf{F} : \mathbb{R}^{n+m} \to \mathbb{R}^{n+m}.
\]
Let $(\mathbf{y}_0,\mathbf{z}_0)$ be given.

Assume there exists an open neighbourhood $\mathcal U \subset \mathbb{R}^{n+m}$
of a solution $\mathbf{x}^* = (\mathbf{y}^*,\mathbf{z}^*)$ such that:

\begin{enumerate}
\item $\mathbf{G},\mathbf{H} \in C^1(\mathcal U)$;

\item the partial Jacobians
\[
\partial_{\mathbf y}\mathbf G(\mathbf y,\mathbf z),
\qquad
\partial_{\mathbf z}\mathbf H(\mathbf y,\mathbf z)
\]
are uniformly nonsingular on $\mathcal U$;

\item there exist constants
$\delta_0^{(z)}, \delta_0^{(y)} > 0$ and
$\delta_{\max}^{(z)}, \delta_{\max}^{(y)} < \infty$
such that the adaptive pseudo-time steps defined below satisfy
\[
0 < \delta_{k,l}^{(z)} \le \delta_{\max}^{(z)},
\qquad
0 < \delta_k^{(y)} \le \delta_{\max}^{(y)};
\]

\item all inner and outer iterates generated by the algorithm remain in
$\mathcal U$.
\end{enumerate}

For each $k \ge 0$, define the inner pseudo-transient continuation iteration
$\{\mathbf{z}_k^{\,l}\}_{l \ge 0}$ by
\[
\left(
(\delta_{k,l}^{(z)})^{-1} \mathbf I
+
\partial_{\mathbf z}\mathbf H(\mathbf y_k, \mathbf z_k^{\,l})
\right)
(\mathbf z_k^{\,l+1} - \mathbf z_k^{\,l})
=
-
\mathbf H(\mathbf y_k, \mathbf z_k^{\,l}),
\]
with adaptive time-step
\[
\delta_{k,l}^{(z)}
=
\min\!\left(
\delta_0^{(z)}
\frac{\|\mathbf H(\mathbf y_k, \mathbf z_k^{\,0})\|}
     {\|\mathbf H(\mathbf y_k, \mathbf z_k^{\,l})\|},
\,
\delta_{\max}^{(z)}
\right).
\]

Define
\[
\mathbf z_k^*
:=
\lim_{l\to\infty} \mathbf z_k^{\,l},
\]
whose existence is guaranteed by pseudo-transient continuation theory.

The outer iteration for $\mathbf y$ is defined by
\[
\left(
(\delta_k^{(y)})^{-1} \mathbf I
+
\partial_{\mathbf y}\mathbf G(\mathbf y_k, \mathbf z_k^*)
\right)
(\mathbf y_{k+1} - \mathbf y_k)
=
-
\mathbf G(\mathbf y_k, \mathbf z_k^*),
\]
with adaptive time-step
\[
\delta_k^{(y)}
=
\min\!\left(
\delta_0^{(y)}
\frac{\|\mathbf F(\mathbf x_0)\|}
     {\|\mathbf F(\mathbf y_k,\mathbf z_k^*)\|},
\,
\delta_{\max}^{(y)}
\right).
\]

Then the sequence $\{(\mathbf y_k,\mathbf z_k^*)\}$ converges to
$\mathbf x^* = (\mathbf y^*,\mathbf z^*)$, and
\[
\mathbf F(\mathbf x^*) = \mathbf 0.
\]
\end{theorem}
\begin{proof}
Fix $k$ and consider the inner problem for $\mathbf z$ with $\mathbf y_k$
held fixed. Define
\[
\mathbf H_{\mathbf y_k}(\mathbf z) := \mathbf H(\mathbf y_k,\mathbf z).
\]

By Assumption~(1), $\mathbf H_{\mathbf y_k} \in C^1$, and by
Assumption~(2), $\partial_{\mathbf z}\mathbf H(\mathbf y_k,\mathbf z)$ is
nonsingular in $\mathcal U$.

The inner iteration is precisely the pseudo-transient continuation method
applied to the nonlinear equation
\[
\mathbf H_{\mathbf y_k}(\mathbf z) = \mathbf 0,
\]
with adaptive pseudo-time step $\delta_{k,l}^{(z)}$ that is uniformly bounded
above and bounded away from zero by Assumption~(3).
Therefore, by the global convergence theorem for pseudo-transient continuation
\cite{kelley1998convergence}, the sequence $\{\mathbf z_k^{\,l}\}_{l\ge0}$
converges to a limit
\[
\mathbf z_k^* := \lim_{l\to\infty} \mathbf z_k^{\,l},
\]
satisfying
\[
\mathbf H(\mathbf y_k,\mathbf z_k^*) = \mathbf 0.
\]

Thus, for each outer iteration index $k$, the inner pseudo-transient
continuation defines a reduced state $\mathbf z_k^*$.

The Implicit Function Theorem guarantees that the reduced state
$\mathbf z_k^*$ depends smoothly on $\mathbf y_k$ in a neighbourhood of
$\mathbf y^*$, and that the solution of
$\mathbf H(\mathbf y,\mathbf z)=\mathbf 0$ is locally unique.
Consequently, the reduced outer problem obtained by fixing
$\mathbf z = \mathbf z_k^*$ is well defined.

Next, define the reduced outer operator
\[
\mathbf G_{\mathbf z_k^*}(\mathbf y) := \mathbf G(\mathbf y,\mathbf z_k^*).
\]
By Assumption~(1), $\mathbf G_{\mathbf z_k^*} \in C^1$, and by
Assumption~(2), $\partial_{\mathbf y}\mathbf G(\mathbf y,\mathbf z_k^*)$ is
nonsingular in a neighbourhood of $\mathbf y^*$.

The outer iteration is again a pseudo-transient continuation method, now
applied to the reduced equation
\[
\mathbf G_{\mathbf z_k^*}(\mathbf y) = \mathbf 0,
\]
with adaptive pseudo-time step $\delta_k^{(y)}$ satisfying the same boundedness
properties. 

Applying the pseudo-transient convergence theorem once more yields
\[
\mathbf y_k \xrightarrow[k\to\infty]{} \mathbf y^*
\quad\text{and}\quad
\mathbf G(\mathbf y^*,\mathbf z^*) = \mathbf 0.
\]

Finally, since $\mathbf H(\mathbf y_k,\mathbf z_k^*) = \mathbf 0$ for all $k$
and $\mathbf H$ is continuous by Assumption~(1), passing to the limit gives
\[
\mathbf H(\mathbf y^*,\mathbf z^*) = \mathbf 0.
\]

Therefore,
\[
\mathbf F(\mathbf x^*) = \mathbf 0,
\]
which completes the proof.
\end{proof}

% --------------------
% Bibliography
% --------------------
\bibliographystyle{elsarticle-num}
\bibliography{convergence}

\end{document}
